\section{Conclusion et perspectives}
\label{sec:conclusion}

\subsection{Bilan}

Le projet livre une plateforme web fonctionnelle qui couvre l'intégralité
du périmètre défini dans la section~\ref{sec:specs-fonc}~: pipeline produit
Kanban, bibliothèque d'angles/créatifs/hooks, suivi budgétaire avec calcul
automatique de ROI, assistant IA multilingue propulsé par DeepSeek, et
modèle prédictif de catégorisation du potentiel d'un nouveau produit. Les
trois langues exigées (anglais, français, arabe) sont supportées
nativement, avec gestion correcte de l'écriture droite-à-gauche pour
l'arabe.

D'un point de vue pédagogique, le projet illustre concrètement~:
\begin{itemize}
    \item La maîtrise du framework Django dans sa configuration moderne
    (templates + HTMX + Tailwind, sans SPA).
    \item L'intégration d'une API LLM tierce dans une application métier,
    avec injection de contexte structuré pour produire des réponses
    \emph{ancrées} dans les données.
    \item La mise en œuvre d'un pipeline ML simple mais complet
    (entraînement, sérialisation, inférence en production).
    \item Le respect des bonnes pratiques d'i18n côté serveur.
\end{itemize}

\subsection{Limites identifiées}

\begin{itemize}
    \item Le jeu de données utilisé pour entraîner le catégoriseur est
    \emph{synthétique}~; les performances en conditions réelles devraient
    être validées avec un historique de campagnes authentiques.
    \item Le chatbot ne dispose pas (encore) de \emph{function-calling}~:
    il peut citer les données du CRM mais ne peut pas, par exemple, créer
    une campagne ou modifier un budget par lui-même.
    \item L'authentification reste basique (Django auth + sessions) ;
    aucune intégration SSO n'est prévue dans cette itération.
\end{itemize}

\subsection{Perspectives d'évolution}

\begin{enumerate}
    \item \textbf{Connecteurs régies publicitaires.} Brancher Meta Ads,
    Google Ads et TikTok Ads via leurs APIs officielles pour automatiser
    la saisie des dépenses et des performances.
    \item \textbf{Génération de créatifs assistée.} Coupler le chatbot à
    un modèle texte-vers-image (Stable Diffusion, FLUX) pour générer des
    propositions de visuels publicitaires directement depuis le CRM.
    \item \textbf{Function-calling DeepSeek.} Faire évoluer le chatbot en
    agent capable d'exécuter des actions (créer une campagne, marquer un
    créatif comme \emph{winner}\ldots) via l'API \emph{tools} de DeepSeek.
    \item \textbf{Prévision time-series.} Ajouter un second modèle
    (Prophet ou LightGBM) pour prévoir les dépenses publicitaires des 30
    prochains jours.
    \item \textbf{Déploiement} containerisé (Docker + docker-compose),
    hébergement sur une PaaS (Fly.io, Railway) avec PostgreSQL managé.
\end{enumerate}

\subsection{Mot de la fin}

Au-delà de la satisfaction des exigences académiques, ce projet a permis
d'expérimenter une démarche de plus en plus répandue dans l'industrie~:
construire des applications métier \textbf{augmentées par l'IA}, où le LLM
n'est plus un simple chatbot décoratif mais un véritable collaborateur
contextualisé. C'est cette intégration profonde, plus que l'usage isolé
d'un modèle de langue, qui crée la valeur perçue par l'utilisateur final.
